% page 537 of the facsimile (image p-536 of the scan)
% block 180.3 x 242.0 mm ; left margin 21.7 mm
\pgc{0.000mm}{0.000mm}{
\l{\cel{58}529}
\l{}
\l{\cel{5}\sou{simetra}. Karakterizata da la korespondo segun-norma (pri grandeso,}
\l{\cel{8}figuro, situeso) inter korpo o parti di korpo. (a) (\sou{anat.})}
\l{\cel{8}Dispozeso di organi dua, situita, ica dextre, ita sinistre, di}
\l{\cel{8}la lineo mediana (che la vertebrozi e la artropodi) ; (b)}
\l{\cel{8}(\sou{geom.}) (Figuro) di qua la elementi reciproka inter-egala esas}
\l{\cel{8}situita inverse. - DEFIRS.}
\l{}
\l{\cel{5}\sou{simfizo}. (\sou{anat.)}\cel{2}Artiko ek du osti, fixa o poke-moviva, per fibro}
\l{\cel{8}kartilago. - DEFIRS.}
\l{}
\l{\cel{5}\sou{simfonio}. (\sou{muziko)}. Muzikajo por plura voco-toni od instrumenti}
\l{\cel{8}koncertala : (a)(\sou{olim}): Muzikajo orkestrala qua esas la uver}
\l{\cel{8}turo di opero ; (b)(de\cel{1}\sou{la yarcento l8-esma}) Kompozajo muzikala}
\l{\cel{8}por orkestro, segun la formulo di sonato, t.e. qua konsistas ye}
\l{\cel{8}prefaco, adajio, skerco e finalo. - DEFIRS.}
\l{}
\l{\cel{5}\sou{simio}. (\sou{zool.}) Mamifero quadrimanua. - L.\cel{1}\sou{simia}. - EFI.}
\l{}
\l{\cel{5}\sou{simila.}\cel{2}Di qua la eso-maniero diferas per nur tre poke de la}
\l{\cel{8}eso-maniero di altra persono, di altra kozo. - EFIS.}
\l{}
\l{\cel{5}\sou{simoniar}. (\sou{netrans}.) Komercachar pri la kozi religiala (sakramen-}
\l{\cel{8}ti, tituli ekleziala, e c.) quin onu donas od aquiras po peku-}
\l{\cel{8}nio o po avantajo sekulara di altra sorto. - DEFIS.}
\l{}
\l{\cel{5}\sou{simpatiar}. (\sou{trans.)}\cel{1}I. Sentar su afina kun altru. - II. Atrakt-}
\l{\cel{8}esar da altru per tendenco instintala. - DEFIRS.}
\l{}
\l{\cel{5}\sou{simpatiko}. (\sou{anat.}) I. Centro nervala aparta de cerebro. - II.}
\l{\cel{8}Nervaro ganglionoza ek dua kordono, dextre e sinistre di la}
\l{\cel{8}spino. -}
\l{}
\l{\cel{5}\sou{simpla}. Qua ne konsistas ek parti. - EFIS.}
\l{}
\l{\cel{5}\sou{simptomo}. I. (\sou{patol.}) Fenomeno qua karakterizas la existesko di}
\l{\cel{8}ta o ca morbo. - II. (\sou{metaf}.) Signo karakteriziva. - DEFIRS.}
\l{}
\l{\cel{5}\sou{simular}. (\sou{trans}.)I. Donar la semblo kom la realajo. - II. (\sou{yuro}-}
\l{\cel{8}\sou{cienco)}. Facar (akto) qua eludas la lego per indiki falsa. -}
\l{\cel{8}DEFIS.}
\l{}
\l{\cel{5}\sou{simultana}. Qua valoras egale e kune, solidare. (Kom ex.: equacioni}
\l{\cel{8}simultana). - \textquotedbl{}Simultana\textquotedbl{} ne esas sinonimo di \textquotedbl{}sama-tempa}
\l{\cel{7}\textquotedbl{}sam-epoka\textquotedbl{}, \textquotedbl{}sam-instanta\textquotedbl{}). -}
\l{}
\l{\cel{5}\sou{sino}. I. (\sou{anat.}) Parto di la korpo di muliero, en qua elu facis}
\l{\cel{8}sua filio. - II. (\sou{metaf.}) Medio. - FIS.}
\l{}
\l{\cel{5}sinagogo. I. (\sou{historio religiala)}. Asemblitaro ek la religiani ,}
\l{\cel{8}ekleziani (lego da Mozes). - II. Religio di la Judi (opoze a}
\l{\cel{8}kristanismo). - III. Templo religiala di la Judi. - DEFIRS.}
\l{}
\l{\cel{5}\sou{sinapo}. (\sou{bot.)}\cel{2}Planto krucifera ; ek lua grani facesas mustardo.}
\l{\cel{8}- L. sinapis. - FIL}
}
